\section{Experimental Validation}
\label{sec:experiments}

\subsection{Objectives}

We answer three validation questions aligned with the event-signal thesis:
\begin{itemize}\itemsep1pt
  \item \textbf{RQ1:} Under strict group isolation, does Stage~1 retain strong OOD macro-F1 suitable for dispatch?
  \item \textbf{RQ2:} On held-out malicious-focused data, does Stage~2 retain non-trivial macro-F1 despite class imbalance?
  \item \textbf{RQ3:} Does replay reproduce one event signal per input flow with consistent Stage~1/Stage~2 admission?
\end{itemize}
Success means downstream consumers can trust serialized fields---not that the pipeline autonomously contains incidents.

\subsection{Data and protocol}

Evaluations fuse controlled laboratory HTTPS/TLS scenarios with public encrypted-traffic corpora (tunneling, DNS-over-HTTPS families, etc.).
Stage~1 estimates benign vs.\ malicious propagation with train/dev/test sizes $(4160,450,460)$, class ratio $3744{:}416$ on training, and \emph{zero} split-group overlap---reducing leakage from correlated captures.
Stage~2 trains on dispatched flows with split $(414,88,93)$; test support per behavior is $(20,14,45,14)$ for C2, data exfiltration, HTTPS tunneling, and domain-fronting-like classes respectively.

All numbers below are \emph{frozen} artifact exports; we do not claim adaptive drift handling or open-set guarantees.

\subsection{Results and interpretation}

Table~\ref{tab:validation} aggregates headline metrics.
Stage~1 achieves \textbf{0.9478} OOD macro-F1, supporting its role as a conservative yet informative gate: analysts receive calibrated context even when flows stay benign, while suspicious flows accumulate structured evidence.

Stage~2 reaches \textbf{0.8495} accuracy and \textbf{0.7778} macro-F1 on $n=93$.
Macro-F1 reflects imbalance; accuracy complements readability but should not be interpreted as operational utility in isolation.
These metrics validate \emph{behavior tagging conditioned on dispatch}, not end-to-end investigative conclusions.

Replay exercises the full serialization path: \textbf{460} ingested flows yield \textbf{460} signals; Stage~1 tallies \textbf{418} benign vs.\ \textbf{42} malicious calls, and all \textbf{42} malicious flows invoke Stage~2---demonstrating parity between scoring, dispatch booleans, and serialized event fields.
Discrepancies here would break SOAR automations; their absence supports the engineering thesis of threshold-consistent events.

\subsection{Limitations}

We do not evaluate long-horizon temporal fusion, earliest-prefix decisions, continual learning under concept drift, or detection of entirely novel threat families.
Nor do we study Edge-hosted large language models or automated blocking policies consuming these signals.
Those directions require consumers and safeguards beyond the scope of this input-layer validation.

\begin{table*}[t]
  \centering
  \caption{Validated experimental snapshot (counts are flows unless noted).}
  \label{tab:validation}
  \footnotesize
  \setlength{\tabcolsep}{4pt}
  \begin{tabular}{@{}p{0.28\textwidth} p{0.66\textwidth}@{}}
    \toprule
    Item & Value \\
    \midrule
    Stage~1 train / dev / test & 4160 / 450 / 460 \\
    Stage~1 train benign:malicious & 3744 : 416 \\
    Stage~1 group overlap (train--dev--test) & 0 \\
    Stage~1 OOD macro-F1 (test) & 0.9478 \\
    \midrule
    Stage~2 train / dev / test & 414 / 88 / 93 \\
    Stage~2 test accuracy / macro-F1 & 0.8495 / 0.7778 \\
    Stage~2 test support by class$^{\dagger}$ & 20 / 14 / 45 / 14 \\
    \midrule
    Replay flows / emitted event signals & 460 / 460 \\
    Replay Stage~1 benign / malicious & 418 / 42 \\
    Replay flows entering Stage~2 & 42 \\
    \bottomrule
  \end{tabular}\\[-2pt]
  {\footnotesize $^{\dagger}$C2, data exfiltration, HTTPS tunneling, domain-fronting-like.}
\end{table*}

